% =============================================================================
% Kapitel 7: Evaluation
% =============================================================================

Dieses Kapitel beschreibt die Evaluation des in Kap.~\ref{sec:implementierung} implementierten Systems.
Zunächst wird die Konstruktion des Gold-Standard-Datensatzes dargestellt, anschließend die Validierung der LLM-gestützten Annotation und abschließend die vergleichende Evaluation der drei Analyseansätze.

\subsection{Gold-Standard-Datensatz}\label{sec:eval-goldstandard}

Die in Kap.~\ref{sec:gold-standard} beschriebene Evaluationsmethodik erfordert einen annotierten Datensatz, der als Ground Truth dient.
Dieser Abschnitt beschreibt die Datenquellen, das Annotationsverfahren und die resultierende Datensatzstruktur.

\subsubsection{Datenquellen}\label{sec:eval-datenquellen}

Um die Generalisierbarkeit der Ergebnisse über eine einzelne Hochschule hinaus zu gewährleisten, wurden Qualifikationsziel-Abschnitte aus Modulhandbüchern von vier Hochschulen und vier Fachbereichen extrahiert.
Tab.~\ref{tab:datenquellen} gibt eine Übersicht.

\begin{table}[H]
\caption{Datenquellen des Gold-Standard-Datensatzes\label{tab:datenquellen}}
\centering
\small
\begin{tabular}{llrr}
\hline
\textbf{Hochschule} & \textbf{Fachbereich} & \textbf{Module} & \textbf{Sätze} \\
\hline
Hochschule Fulda & Angewandte Informatik & 391 & 3\,213 \\
TU Darmstadt & Informatik & 41 & 896 \\
TH Mittelhessen & Maschinenbau & 73 & 281 \\
TH Mittelhessen & Wirtschaftsingenieurwesen & 63 & 1\,187 \\
Uni Kassel & Soziale Arbeit & 17 & 40 \\
\hline
\textbf{Gesamt} & \textbf{4 Fachbereiche} & \textbf{585} & \textbf{5\,617} \\
\hline
\end{tabular}
\end{table}

Die Auswahl der externen Hochschulen folgt drei Kriterien: (1)~öffentliche Verfügbarkeit der Modulhandbücher als PDF, (2)~Diversität der Fachbereiche, um domänenübergreifende Robustheit zu prüfen, und (3)~unterschiedliche Formulierungstraditionen, da ingenieurwissenschaftliche und sozialwissenschaftliche Modulhandbücher andere sprachliche Muster aufweisen als informatische.

Die Textextraktion erfolgt mit \texttt{pdftotext} (Poppler-Bibliothek), da komplexe PDF-Layouts mit Python-Bibliotheken wie pdfplumber zu Abstürzen führten.
Für jede Hochschule wurde ein spezifischer Extraktor implementiert, der die jeweilige Dokumentstruktur berücksichtigt.
Kurze Fragmente ($< 10$~Zeichen) sowie Abschnittsüberschriften werden per Regex-Filter entfernt.

\subsubsection{Annotationsschema}\label{sec:eval-schema}

Jeder Satz wird in drei Dimensionen annotiert:

\begin{enumerate}
  \item \textbf{Satztyp} (\texttt{typ}): Kompetenzformulierung~(K), Inhaltsbeschreibung~(I) oder Sonstiges~(S), gemäß den Entscheidungsregeln in Kap.~\ref{sec:eval-methodik}.
  \item \textbf{Taxonomiestufe} (\texttt{taxonomie}): Für Kompetenzformulierungen die kognitive Prozessstufe nach Anderson und Krathwohl~\cite{anderson2001taxonomy} (Stufen 1--6). Ein Satz kann mehreren Stufen zugeordnet werden, wenn er mehrere Handlungsverben auf verschiedenen Stufen enthält (z.\,B. \enquote{analysieren und anwenden} $\rightarrow$ Stufen~3;4).
  \item \textbf{Formulierungsqualität} (\texttt{qualitaet}): Eine dreistufige Bewertung (gut, akzeptabel, schlecht), die erfasst, ob die Formulierung den Kriterien für outcome-orientierte Kompetenzformulierungen entspricht~\cite{hrknexus2015lernergebnisse}.
\end{enumerate}

Die Mehrfachzuordnung bei der Taxonomie unterscheidet den Datensatz von bestehenden Ansätzen, die typischerweise die höchste Stufe als alleinige Annotation verwenden.
Dies ermöglicht eine differenziertere Evaluation, da das System auch Nebenverben korrekt erkennen muss.

\subsubsection{LLM-gestützte Vorannotation}\label{sec:eval-vorannotation}

Die manuelle Annotation von 5\,617~Sätzen in drei Dimensionen wäre mit einem Aufwand von geschätzt 100--150~Stunden verbunden.
Um diesen Aufwand zu reduzieren und gleichzeitig eine konsistente Ausgangsbasis zu schaffen, wird ein zweistufiges Annotationsverfahren eingesetzt.

\paragraph{Stufe 1: LLM-Vorannotation.}
Alle Sätze werden automatisiert durch Claude Sonnet annotiert.\footnote{Modell: \texttt{claude-sonnet-4-6}, Anthropic API, Temperatur~0.}
Das LLM erhält den vollständigen Annotationsleitfaden als Systemprompt und verarbeitet die Sätze in Batches von je 25~Sätzen.
Neben den drei Annotationsdimensionen gibt das Modell für jeden Satz einen Konfidenzwert zwischen 0.0 und 1.0 aus, der die Selbsteinschätzung der Annotationssicherheit widerspiegelt.

Aktuelle Forschung zeigt, dass LLM-Annotationen bei verschiedenen Klassifikationsaufgaben mit menschlichen Annotatoren vergleichbar oder überlegen sein können.
Gilardi et al.~\cite{gilardi2023chatgpt} berichten, dass ChatGPT bei vier Textklassifikations-Tasks Crowd-Worker um durchschnittlich 25~Prozentpunkte in der Accuracy übertrifft, bei Kosten von unter \$0,003 pro Annotation.
Törnberg~\cite{tornberg2025llm} bestätigt dies für politische Textklassifikation: GPT-4 erreicht höhere Accuracy und Reliabilität als sowohl Crowd-Worker als auch trainierte Experten.

Gleichzeitig warnen Pangakis et al.~\cite{pangakis2023validation}, dass die Performance stark je nach Aufgabentyp variiert -- bei 9 von 27~replizierten Tasks lag Precision oder Recall unter~0,5, trotz hoher Gesamt-Accuracy.
Reiss~\cite{reiss2023testing} weist zudem auf die Nicht-Determiniertheit von LLMs hin: Selbst bei identischem Input können verschiedene Durchläufe unterschiedliche Ergebnisse erzeugen.
Törnberg~\cite{tornberg2024bestpractices} empfiehlt daher explizit, LLMs in einen rigorosen Annotationsprozess einzubetten, der systematische Validierung gegen menschliche Labels umfasst.

Das hier eingesetzte zweistufige Verfahren -- LLM-Vorannotation mit anschließender stratifizierter manueller Validierung und Cohen's~$\kappa$ -- entspricht diesen Best Practices.

\paragraph{Stufe 2: Manuelle Validierung.}
Eine stratifizierte Stichprobe von 300~Sätzen wird manuell überprüft und bei Bedarf korrigiert.
Die Stratifizierung erfolgt nach drei Kriterien: Hochschule, Satztyp und LLM-Konfidenz (unterteilt in niedrig $< 0.7$, mittel $0.7$--$0.9$, hoch $> 0.9$).
Dadurch werden systematisch auch unsichere Annotationen und unterrepräsentierte Hochschulen in der Validierung berücksichtigt.

\subsubsection{Inter-Annotator-Reliabilität}\label{sec:eval-kappa}

Um die Qualität der Annotationen unabhängig von der LLM-Vorannotation zu bewerten, wird die Inter-Annotator-Reliabilität (IAR) bestimmt.
Ein zweiter, unabhängiger Annotator annotiert dieselben 300~Sätze der Validierungsstichprobe, ohne Kenntnis der LLM- oder Erstannotation.

Die Übereinstimmung wird mittels Cohen's $\kappa$~\cite{cohen1960coefficient} berechnet, das die beobachtete Übereinstimmung um die zufallsbedingte Übereinstimmung korrigiert:

\begin{equation}\label{eq:kappa}
\kappa = \frac{p_o - p_e}{1 - p_e}
\end{equation}

\noindent wobei $p_o$ die beobachtete und $p_e$ die erwartete Übereinstimmung bei Zufall bezeichnet.
Die Bewertung folgt der Skala von Landis und Koch~\cite{landis1977measurement}: $\kappa < 0.20$~(gering), $0.21$--$0.40$~(ausreichend), $0.41$--$0.60$~(moderat), $0.61$--$0.80$~(substanziell), $0.81$--$1.00$~(fast perfekt).

Tab.~\ref{tab:kappa} zeigt die Ergebnisse für die drei Annotationsdimensionen.

\begin{table}[H]
\caption{Inter-Annotator-Reliabilität (Cohen's $\kappa$)\label{tab:kappa}}
\centering
\begin{tabular}{llrrrl}
\hline
\textbf{Dimension} & \textbf{Klassen} & \textbf{Erst\,$\leftrightarrow$\,LLM} & \textbf{Zweit\,$\leftrightarrow$\,LLM} & \textbf{Erst\,$\leftrightarrow$\,Zweit} & \textbf{Bewertung} \\
\hline
Satztyp (K/I/S) & 3 & 0{,}623 & 0{,}768 & 0{,}737 & substanziell \\
Taxonomie (Primärstufe) & 6 & 0{,}852 & 0{,}806 & 0{,}861 & fast perfekt \\
Qualität (gut/akz./schl.) & 3 & 0{,}350 & 0{,}401 & 0{,}869 & siehe Text \\
\hline
\end{tabular}
\end{table}

Die Ergebnisse zeigen ein differenziertes Bild.
Die \textbf{Taxonomie-Zuordnung} erreicht mit $\kappa \geq 0{,}806$ in allen drei Paaren die höchste Übereinstimmung -- die Primärstufe nach Anderson und Krathwohl lässt sich offenbar auch ohne fachwissenschaftlichen Hintergrund reliabel annotieren.
Zwischen Erstannotator und Zweitannotator liegt $\kappa = 0{,}861$ (fast perfekt), was die Konsistenz des Annotationsleitfadens bestätigt.

Beim \textbf{Satztyp} liegt die Übereinstimmung zwischen Zweitannotator und LLM ($\kappa = 0{,}768$) höher als zwischen Erstannotator und LLM ($\kappa = 0{,}623$).
Die Inter-Annotator-Übereinstimmung ($\kappa = 0{,}737$) ist substanziell.
Abweichungen betreffen vor allem die Abgrenzung zwischen Kompetenzformulierungen und Sonstigem bei Infinitivkonstruktionen ohne explizites Subjekt.

Die \textbf{Qualitätsbewertung} zeigt ein auffälliges Muster: Während die beiden menschlichen Annotatoren nahezu perfekt übereinstimmen ($\kappa = 0{,}869$), weichen beide deutlich vom LLM ab ($\kappa = 0{,}350$ bzw.\ $0{,}401$).
Das LLM vergab 19-mal die Bewertung \enquote{schlecht}, die menschlichen Annotatoren hingegen nur 0- bzw.\ 2-mal.
Diese systematische Abweichung deutet darauf hin, dass das LLM einen strengeren Qualitätsmaßstab anlegt als die menschlichen Annotatoren.
Da die Inter-Annotator-Reliabilität für die Qualitätsdimension hoch ist, wird für die Evaluation der menschliche Konsens als Referenz verwendet.

\subsubsection{Datensatzstatistiken}\label{sec:eval-stats}

Der finale Gold-Standard umfasst 5\,617~Sätze, die stratifiziert nach Hochschule und Satztyp in ein Trainings-Set (80\,\%, $n = 4\,493$) und ein Test-Set (20\,\%, $n = 1\,124$) aufgeteilt werden (Seed~42).
Tab.~\ref{tab:datensatz-stats} zeigt die Verteilung nach Hochschule und Satztyp.

\begin{table}[H]
\caption{Verteilung des Gold-Standard-Datensatzes nach Satztyp\label{tab:datensatz-stats}}
\centering
\small
\begin{tabular}{lrrrr}
\hline
\textbf{Hochschule} & \textbf{K} & \textbf{I} & \textbf{S} & \textbf{Gesamt} \\
\hline
Hochschule Fulda & 3\,049 & 66 & 98 & 3\,213 \\
TU Darmstadt & 684 & 104 & 108 & 896 \\
TH Mittelhessen & 1\,404 & 18 & 46 & 1\,468 \\
Uni Kassel & 37 & 0 & 3 & 40 \\
\hline
\textbf{Gesamt} & \textbf{5\,174} & \textbf{188} & \textbf{255} & \textbf{5\,617} \\
\hline
\end{tabular}
\end{table}

Der Datensatz ist erwartungsgemäß unbalanciert: 92{,}1\,\% der Sätze sind Kompetenzformulierungen~(K), 3{,}3\,\% Inhaltsbeschreibungen~(I) und 4{,}5\,\% Sonstiges~(S).
TU Darmstadt weist mit 23{,}7\,\% den höchsten Anteil an Nicht-Kompetenz-Sätzen (I+S) auf, während Hochschule Fulda mit 5{,}1\,\% den niedrigsten hat.

31{,}5\,\% der Kompetenzformulierungen sind mehreren Taxonomiestufen zugeordnet.
Tab.~\ref{tab:taxonomie-verteilung} zeigt die Verteilung der Taxonomiestufen über den gesamten Datensatz.

\begin{table}[H]
\caption{Taxonomiestufen-Verteilung (Mehrfachzuordnung möglich)\label{tab:taxonomie-verteilung}}
\centering
\begin{tabular}{lrr}
\hline
\textbf{Stufe} & \textbf{Anzahl} & \textbf{Anteil} \\
\hline
1 -- Erinnern & 1\,082 & 14{,}9\,\% \\
2 -- Verstehen & 1\,355 & 18{,}6\,\% \\
3 -- Anwenden & 2\,376 & 32{,}7\,\% \\
4 -- Analysieren & 682 & 9{,}4\,\% \\
5 -- Bewerten & 1\,027 & 14{,}1\,\% \\
6 -- Erschaffen & 749 & 10{,}3\,\% \\
\hline
\end{tabular}
\end{table}

Stufe~3 (Anwenden) dominiert mit 32{,}7\,\%, während die höheren Stufen~4--6 zusammen 33{,}8\,\% ausmachen.
Die Qualitätsverteilung zeigt 35{,}2\,\% \enquote{gut}, 52{,}3\,\% \enquote{akzeptabel} und 12{,}4\,\% \enquote{schlecht}.

% =============================================================================
\subsection{Vergleichende Evaluation}\label{sec:eval-vergleich}

Die drei in Kap.~\ref{sec:entwurf} definierten Analyseansätze -- regelbasiert, Embedding-basiert und hybrid -- werden auf dem Test-Set des Gold-Standards evaluiert.
Für jeden Ansatz werden die in Kap.~\ref{sec:metriken} definierten Metriken berechnet.

\subsubsection{Klassifikation}\label{sec:eval-klassifikation}

Tab.~\ref{tab:eval-klassifikation} zeigt die Ergebnisse der Satztyp-Klassifikation (K/I/S) für die drei Ansätze.

\begin{table}[H]
\caption{Klassifikationsergebnisse (Satztyp K/I/S)\label{tab:eval-klassifikation}}
\centering
\begin{tabular}{lrrr}
\hline
\textbf{Ansatz} & \textbf{Precision} & \textbf{Recall} & \textbf{F1 (weighted)} \\
\hline
Regelbasiert & 0{,}899 & 0{,}767 & 0{,}819 \\
Embedding (SVM) & 0{,}964 & 0{,}966 & 0{,}963 \\
Hybrid & 0{,}945 & 0{,}936 & 0{,}929 \\
\hline
\end{tabular}
\end{table}

Der Embedding-Ansatz erreicht mit $F_1 = 0{,}963$ das beste Ergebnis und übertrifft den regelbasierten Ansatz ($F_1 = 0{,}819$) deutlich.
Der hybride Ansatz liegt mit $F_1 = 0{,}929$ dazwischen.
Die gewichteten Metriken werden von der dominanten K-Klasse (92{,}1\,\%) bestimmt; die klassenweisen Ergebnisse in Tab.~\ref{tab:eval-klassifikation-detail} zeigen ein differenzierteres Bild.

\begin{table}[H]
\caption{Klassenweise F1-Scores der Satztyp-Klassifikation\label{tab:eval-klassifikation-detail}}
\centering
\begin{tabular}{lrrr}
\hline
\textbf{Klasse} & \textbf{Regelbasiert} & \textbf{Embedding} & \textbf{Hybrid} \\
\hline
K (Kompetenz) & 0{,}881 & 0{,}984 & 0{,}976 \\
I (Inhalt) & 0{,}000 & 0{,}645 & 0{,}233 \\
S (Sonstiges) & 0{,}177 & 0{,}778 & 0{,}487 \\
\hline
\end{tabular}
\end{table}

Der regelbasierte Ansatz erkennt Inhaltsbeschreibungen überhaupt nicht ($F_1 = 0{,}000$), da er ohne Verbliste-Treffer keine Unterscheidung zwischen I und S vornehmen kann.
Der Embedding-Ansatz erzielt hier $F_1 = 0{,}645$ (I) bzw.\ $F_1 = 0{,}778$ (S) -- die semantische Repräsentation ermöglicht eine Unterscheidung, die rein lexikalisch nicht möglich ist.
Der hybride Ansatz erreicht bei I zwar perfekte Precision ($1{,}000$), aber sehr niedrigen Recall ($0{,}132$): Er erkennt I-Sätze nur, wenn das Embedding eindeutig ist, klassifiziert aber viele als K.

\subsubsection{Taxonomie-Zuordnung}\label{sec:eval-taxonomie}

Tab.~\ref{tab:eval-taxonomie} zeigt die Ergebnisse der Taxonomie-Zuordnung für Kompetenzformulierungen.
Da der regelbasierte Ansatz nur Sätze mit Verbliste-Treffern klassifizieren kann, werden 536 von 1\,035~K-Sätzen (51{,}8\,\%) nicht zugeordnet.
Die Metriken beziehen sich jeweils nur auf die klassifizierten Sätze.

\begin{table}[H]
\caption{Taxonomie-Zuordnung (Primärstufe 1--6, nur K-Sätze)\label{tab:eval-taxonomie}}
\centering
\begin{tabular}{lrrrr}
\hline
\textbf{Ansatz} & \textbf{Precision} & \textbf{Recall} & \textbf{F1 (weighted)} & \textbf{Abdeckung} \\
\hline
Regelbasiert & 0{,}639 & 0{,}585 & 0{,}581 & 48{,}2\,\% \\
Embedding (SVM) & 0{,}842 & 0{,}842 & 0{,}839 & 100\,\% \\
Hybrid & 0{,}739 & 0{,}706 & 0{,}713 & 89{,}0\,\% \\
\hline
\end{tabular}
\end{table}

Der Embedding-Ansatz klassifiziert alle K-Sätze und erreicht $F_1 = 0{,}839$.
Der regelbasierte Ansatz erreicht auf den 48{,}2\,\% der Sätze, die er abdeckt, nur $F_1 = 0{,}581$.
Der hybride Ansatz verbessert die Abdeckung auf 89{,}0\,\%, bleibt aber mit $F_1 = 0{,}713$ hinter dem reinen Embedding-Ansatz zurück.

Tab.~\ref{tab:eval-taxonomie-detail} schlüsselt die Ergebnisse nach Taxonomiestufe auf.

\begin{table}[H]
\caption{F1-Score nach Taxonomiestufe\label{tab:eval-taxonomie-detail}}
\centering
\begin{tabular}{lrrr}
\hline
\textbf{Stufe} & \textbf{Regelbasiert} & \textbf{Embedding} & \textbf{Hybrid} \\
\hline
1 -- Erinnern & 0{,}129 & 0{,}860 & 0{,}664 \\
2 -- Verstehen & 0{,}735 & 0{,}849 & 0{,}753 \\
3 -- Anwenden & 0{,}665 & 0{,}870 & 0{,}791 \\
4 -- Analysieren & 0{,}305 & 0{,}667 & 0{,}450 \\
5 -- Bewerten & 0{,}577 & 0{,}806 & 0{,}649 \\
6 -- Erschaffen & 0{,}533 & 0{,}766 & 0{,}567 \\
\hline
\end{tabular}
\end{table}

Der Embedding-Ansatz übertrifft den regelbasierten in jeder Stufe.
Besonders auffällig ist Stufe~1 (Erinnern): Der regelbasierte Ansatz erreicht nur $F_1 = 0{,}129$, da viele Wissensverben (z.\,B.\ \enquote{kennen}, \enquote{wissen}) nicht in der statischen Verbliste enthalten sind oder mit anderen Stufen geteilt werden.
Stufe~4 (Analysieren) ist für alle Ansätze die schwierigste, was mit der geringen Trainingsdatenmenge (9{,}4\,\% des Datensatzes) und der semantischen Nähe zu Stufe~3 zusammenhängt.
Stufe~3 (Anwenden) erreicht beim Embedding-Ansatz mit $F_1 = 0{,}870$ den besten Wert, begünstigt durch den größten Stichprobenumfang (32{,}7\,\%).

\subsubsection{Aufschlüsselung nach Schwierigkeit}\label{sec:eval-schwierigkeit}

Um die Stärken und Schwächen der Ansätze differenziert zu bewerten, werden die Ergebnisse nach drei Schwierigkeitskategorien aufgeschlüsselt:

\begin{enumerate}
  \item \textbf{Standard}: Sätze mit Verben, die eindeutig in der Verbliste enthalten sind und deren Taxonomiestufe unstrittig ist.
  \item \textbf{Unbekannte Verben}: Kompetenzformulierungen mit Verben, die nicht in der statischen Verbliste des regelbasierten Systems vorkommen (z.\,B. Synonyme, fachspezifische Verben).
  \item \textbf{Nicht-Kompetenz}: Inhaltsbeschreibungen und Sonstiges -- Sätze, die das System korrekt als Nicht-Kompetenz erkennen muss.
\end{enumerate}

\begin{table}[H]
\caption{F1-Score nach Schwierigkeitskategorie\label{tab:eval-schwierigkeit}}
\centering
\begin{tabular}{lrrr}
\hline
\textbf{Kategorie} & \textbf{Regelbasiert} & \textbf{Embedding} & \textbf{Hybrid} \\
\hline
Standard & 0{,}994 & 0{,}999 & 1{,}000 \\
Unbekannte Verben & 0{,}683 & 0{,}996 & 0{,}976 \\
Nicht-Kompetenz & 0{,}279 & 0{,}740 & 0{,}425 \\
\hline
\end{tabular}
\end{table}

Bei Standardfällen -- Sätze mit eindeutigen Verbliste-Treffern -- erreichen alle drei Ansätze nahezu perfekte Klassifikation ($F_1 \geq 0{,}994$).
Die Unterschiede zeigen sich bei den schwierigeren Kategorien: Bei unbekannten Verben bricht der regelbasierte Ansatz auf $F_1 = 0{,}683$ ein, da er Kompetenzformulierungen mit Verben außerhalb der statischen Liste nicht zuverlässig erkennt.
Der Embedding-Ansatz bleibt mit $F_1 = 0{,}996$ stabil, der hybride Ansatz mit $F_1 = 0{,}976$ nur geringfügig darunter.
Bei Nicht-Kompetenz-Sätzen ist der Abstand am größten: Der regelbasierte Ansatz erreicht nur $F_1 = 0{,}279$, da er ohne Verbliste-Evidenz kaum zwischen Inhaltsbeschreibungen und Kompetenzformulierungen differenzieren kann.

\subsubsection{Aufschlüsselung nach Hochschule}\label{sec:eval-hochschule}

Um die Generalisierbarkeit zu prüfen, werden die Ergebnisse nach Hochschule aufgeschlüsselt.
Da das System primär auf Daten der Hochschule Fulda trainiert wurde, ist ein Leistungsabfall bei externen Hochschulen zu erwarten.

\begin{table}[H]
\caption{F1-Score (Klassifikation) nach Hochschule\label{tab:eval-hochschule}}
\centering
\begin{tabular}{lrrr}
\hline
\textbf{Hochschule} & \textbf{Regelbasiert} & \textbf{Embedding} & \textbf{Hybrid} \\
\hline
HS Fulda (Informatik) & 0{,}874 & 0{,}985 & 0{,}955 \\
TU Darmstadt (Informatik) & 0{,}631 & 0{,}867 & 0{,}759 \\
TH Mittelhessen (Ing.) & 0{,}810 & 0{,}973 & 0{,}969 \\
Uni Kassel (Soz. Arbeit)$^*$ & 0{,}923 & 1{,}000 & 1{,}000 \\
\hline
\multicolumn{4}{l}{\scriptsize $^*$\,$n = 7$ im Test-Set; F1-Werte eingeschränkt interpretierbar.} \\
\end{tabular}
\end{table}

Der Embedding-Ansatz generalisiert deutlich besser über Hochschulen hinweg als der regelbasierte.
Hochschule Fulda als Hauptdatenquelle zeigt erwartungsgemäß die besten Ergebnisse ($F_1 = 0{,}985$).
TU Darmstadt weist den stärksten Leistungsabfall auf ($F_1 = 0{,}867$), was mit dem hohen Anteil an Nicht-Kompetenz-Sätzen (23{,}7\,\%) zusammenhängt, die in den informatischen Modulhandbüchern der TU Darmstadt häufiger auftreten.
TH Mittelhessen erreicht trotz des ingenieurwissenschaftlichen Fachvokabulars $F_1 = 0{,}973$, was die domänenübergreifende Robustheit des Embedding-Ansatzes bestätigt.
Uni Kassel (Soziale Arbeit) erreicht $F_1 = 1{,}000$ bei beiden nicht-regelbasierten Ansätzen, wobei die geringe Stichprobengröße ($n = 7$ im Test-Set) die Aussagekraft einschränkt.

Ergänzend zeigt Tab.~\ref{tab:eval-hochschule-tax} die Taxonomie-Zuordnung nach Hochschule.

\begin{table}[H]
\caption{F1-Score (Taxonomie) nach Hochschule\label{tab:eval-hochschule-tax}}
\centering
\begin{tabular}{lrrr}
\hline
\textbf{Hochschule} & \textbf{Regelbasiert} & \textbf{Embedding} & \textbf{Hybrid} \\
\hline
HS Fulda (Informatik) & 0{,}568 & 0{,}911 & 0{,}755 \\
TU Darmstadt (Informatik) & 0{,}680 & 0{,}770 & 0{,}743 \\
TH Mittelhessen (Ing.) & 0{,}564 & 0{,}717 & 0{,}624 \\
Uni Kassel (Soz. Arbeit)$^*$ & 0{,}444 & 0{,}714 & 0{,}381 \\
\hline
\multicolumn{4}{l}{\scriptsize $^*$\,$n = 7$ im Test-Set; F1-Werte eingeschränkt interpretierbar.} \\
\end{tabular}
\end{table}

Bei der Taxonomie-Zuordnung zeigt sich ein ähnliches Muster: Der Embedding-Ansatz erreicht durchgehend die höchsten F1-Werte.
Der regelbasierte Ansatz erzielt bei TU Darmstadt ($F_1 = 0{,}680$) einen höheren Wert als bei HS Fulda ($F_1 = 0{,}568$), da er dort auf weniger, aber eindeutigere Verbliste-Treffer trifft.
TH Mittelhessen zeigt bei allen Ansätzen niedrigere Taxonomie-Werte, was auf das ingenieurwissenschaftliche Fachvokabular zurückzuführen ist, das sich von den informatischen Trainingsverben unterscheidet.

% =============================================================================
\subsection{Empfehlungssystem}\label{sec:eval-empfehlung}

Die Evaluation des Empfehlungssystems (FA3) erfolgt mittels Precision@k: Unter den Top-$k$ vorgeschlagenen Referenzformulierungen wird der Anteil relevanter Ergebnisse gemessen.
Die Relevanz wird durch Übereinstimmung der Taxonomiestufe operationalisiert.

\begin{table}[H]
\caption{Precision@k des Empfehlungssystems (SBERT-Ähnlichkeitssuche)\label{tab:eval-precision-k}}
\centering
\begin{tabular}{lr}
\hline
\textbf{Metrik} & \textbf{Precision} \\
\hline
Precision@1 & 0{,}510 \\
Precision@3 & 0{,}479 \\
Precision@5 & 0{,}465 \\
\hline
\end{tabular}
\end{table}

Das Empfehlungssystem verwendet unabhängig vom Klassifikationsansatz dieselbe SBERT-basierte Ähnlichkeitssuche über die Referenzdatenbank (4\,918~Sätze nach Qualitätsfilterung).
Die Relevanz wird über die Übereinstimmung der Taxonomiestufe operationalisiert: Eine Empfehlung gilt als relevant, wenn ihre Primärstufe mit der des Eingabesatzes übereinstimmt.

Precision@1 $= 0{,}510$ bedeutet, dass die ähnlichste Referenzformulierung in 51\,\% der Fälle dieselbe Taxonomiestufe aufweist.
Zum Vergleich: Die häufigste Taxonomiestufe (Stufe~3, Anwenden) umfasst 32{,}7\,\% des Datensatzes; eine naive Baseline, die stets Formulierungen dieser Mehrheitsklasse empfehlen würde, erreichte $P@1 \approx 0{,}327$.
Die erzielten Werte liegen somit deutlich über dieser Baseline.
Der moderate Abfall auf Precision@5 $= 0{,}465$ zeigt, dass auch bei breiteren Empfehlungslisten knapp die Hälfte der Vorschläge taxonomisch passend ist.
Diese Werte sind im Kontext zu interpretieren: Semantische Ähnlichkeit korreliert mit, ist aber nicht identisch zu taxonomischer Übereinstimmung -- ein Satz auf Stufe~3 (Anwenden) kann semantisch einem Satz auf Stufe~4 (Analysieren) ähneln, wenn beide thematisch verwandt sind.

% =============================================================================
\subsection{Set-Analyse}\label{sec:eval-set}

Die Set-Analyse (FA4) wird qualitativ evaluiert, da keine quantitative Ground Truth auf Modulebene vorliegt.
Anhand ausgewählter Module wird geprüft, ob die automatisch ermittelte Taxonomiestufen-Verteilung und die daraus abgeleiteten Empfehlungen fachlich plausibel sind.

Die Set-Analyse aggregiert die Einzelsatz-Ergebnisse auf Modulebene und visualisiert die Taxonomiestufen-Verteilung als Balkendiagramm.
Dadurch wird sichtbar, ob ein Modul alle kognitiven Stufen abdeckt oder sich auf wenige konzentriert.

Als Beispiel sei ein Informatik-Modul der Hochschule Fulda mit sechs Qualifikationszielen betrachtet.
Das System erkennt vier Sätze auf Stufe~3 (Anwenden) und je einen auf Stufe~2 (Verstehen) und Stufe~5 (Bewerten).
Die Stufen~1, 4 und 6 sind nicht vertreten -- das System generiert daraufhin den Hinweis, dass höhere kognitive Stufen unterrepräsentiert sind.
Eine solche Rückmeldung entspricht der didaktischen Empfehlung, Qualifikationsziele über mehrere Taxonomiestufen zu streuen~\cite{anderson2001taxonomy}.

Die Plausibilität der Set-Analyse hängt direkt von der Qualität der Einzelsatz-Klassifikation ab.
Da der Embedding-Ansatz mit $F_1 = 0{,}839$ die zuverlässigste Taxonomie-Zuordnung liefert, ist die darauf basierende Set-Analyse am aussagekräftigsten.
Der regelbasierte Ansatz kann aufgrund seiner geringen Abdeckung (48{,}2\,\%) kein vollständiges Bild der Taxonomie-Verteilung eines Moduls liefern.

% =============================================================================
\subsection{Konfidenz-Analyse}\label{sec:eval-konfidenz}

Um die Nutzbarkeit der Konfidenzwerte zu bewerten, werden die Klassifikationsergebnisse des Embedding-Ansatzes nach Konfidenzbereichen aufgeschlüsselt.
Die Konfidenz ergibt sich aus den SVM-Wahrscheinlichkeitsschätzungen (\texttt{predict\_proba}).
Tab.~\ref{tab:eval-konfidenz} zeigt die Ergebnisse.

\begin{table}[H]
\caption{F1-Score des Embedding-Ansatzes nach Konfidenzbereich\label{tab:eval-konfidenz}}
\centering
\begin{tabular}{lrrr}
\hline
\textbf{Konfidenz} & \textbf{Anteil} & \textbf{Typ-F1} & \textbf{Tax-F1} \\
\hline
Hoch ($\geq 0{,}8$) & 94{,}8\,\% & 0{,}986 & 0{,}918 \\
Mittel ($0{,}5$--$0{,}8$) & 4{,}7\,\% & 0{,}585 & 0{,}694 \\
Niedrig ($< 0{,}5$) & 0{,}5\,\% & 0{,}250 & 0{,}434 \\
\hline
\end{tabular}
\end{table}

Die Ergebnisse zeigen, dass die Konfidenzwerte eine sinnvolle Qualitätsindikation liefern: Im Hoch-Konfidenz-Bereich ($\geq 0{,}8$) steigt der Typ-F1 auf $0{,}986$ -- ein Zugewinn von $2{,}3$ Prozentpunkten gegenüber der Gesamtpopulation.
Bei der Taxonomie ist der Effekt noch deutlicher: $F_1 = 0{,}918$ gegenüber $0{,}839$ gesamt, bei einer Abdeckung von 73{,}8\,\%.
Korrekt klassifizierte Sätze weisen eine durchschnittliche Typ-Konfidenz von $0{,}971$ auf, fehlklassifizierte nur $0{,}712$.
Bei der Taxonomie beträgt die mittlere Konfidenz $0{,}877$ (korrekt) gegenüber $0{,}676$ (falsch).
Die Konfidenzwerte sind somit ein brauchbarer -- wenn auch unkalibrierter -- Indikator für die Vorhersagezuverlässigkeit.

Die Konfidenzwerte werden im Frontend pro Satz angezeigt und farblich kodiert (grün $\geq 80\,\%$, gelb $50$--$80\,\%$, rot $< 50\,\%$).
Dies ermöglicht es Lehrenden, unsichere Zuordnungen gezielt zu überprüfen.

% =============================================================================
\subsection{Qualitative Fehleranalyse}\label{sec:eval-fehleranalyse}

Die aggregierten Metriken zeigen \emph{wie oft} die Ansätze korrekt klassifizieren, lassen aber offen, \emph{warum} Fehler auftreten.
Dieser Abschnitt illustriert anhand konkreter Beispielsätze aus dem Test-Set typische Fehlermuster.
Tab.~\ref{tab:eval-beispiele} zeigt je ein Beispiel für fünf charakteristische Kategorien.

\begin{table}[H]
\caption{Illustrative Beispielsätze für typische Fehlermuster\label{tab:eval-beispiele}}
\centering
\small
\begin{tabular}{p{0.38\textwidth}p{0.12\textwidth}p{0.12\textwidth}p{0.12\textwidth}p{0.12\textwidth}}
\hline
\textbf{Satz (gekürzt)} & \textbf{Gold} & \textbf{Regel} & \textbf{Embed.} & \textbf{Hybrid} \\
\hline
\enquote{Die Studierenden automatisieren und orchestrieren Data Science Workflows.}
  & K / St.\,3 & S / -- & K / St.\,3 & S / -- \\[4pt]
\enquote{Die Studierenden entwickeln ein Bewusstsein für die Rolle des Fachgebiets MCI.}
  & K / St.\,2 & K / St.\,6 & K / St.\,2 & K / St.\,6 \\[4pt]
\enquote{Die Studierenden erkennen verschiedene algebraische Strukturen.}
  & K / St.\,4 & -- & K / St.\,1 & K / St.\,1 \\[4pt]
\enquote{Sie erhalten Anleitung in der systematischen Planung [\ldots] des Arbeitsprozesses.}
  & S & S & K & K \\[4pt]
\enquote{Die Thematik wird im Team bearbeitet und das Ergebnis präsentiert und diskutiert.}
  & S & S & K & S \\
\hline
\end{tabular}
\end{table}

Die Beispiele verdeutlichen drei wiederkehrende Muster:

\paragraph{Unbekannte Verben.}
\enquote{Automatisieren} und \enquote{orchestrieren} fehlen in der statischen Verbliste.
Der regelbasierte Ansatz kann den Satz nicht als K erkennen; der Embedding-Ansatz erfasst die semantische Handlungsorientierung und klassifiziert korrekt.
Dieses Muster erklärt den Großteil der 235~Fälle im Test-Set, in denen der Embedding-Ansatz beim Satztyp korrekt ist und der regelbasierte nicht.
Im hybriden Cascading wird die fehlerhafte regelbasierte Entscheidung nicht korrigiert, da das Embedding-Fallback erst bei fehlender Regelentscheidung aktiviert wird.

\paragraph{Kontextabhängige Verben.}
\enquote{Entwickeln} ist in der Verbliste als Stufe~6 (Erschaffen) hinterlegt, meint im Kontext \enquote{ein Bewusstsein entwickeln} jedoch Stufe~2 (Verstehen).
Ebenso ist \enquote{erkennen} in der Verbliste Stufe~4 (Analysieren) zugeordnet, beschreibt im Kontext \enquote{algebraische Strukturen erkennen} aber eher Stufe~1 (Erinnern) oder Stufe~4.
Der Embedding-Ansatz nutzt den Satzkontext und ordnet den ersten Fall korrekt als Stufe~2 ein; beim zweiten Beispiel scheitert auch er -- Stufe~4 (Analysieren) ist mit $F_1 = 0{,}667$ die schwächste Kategorie.
Von 82~Stufe-4-Sätzen im Test-Set werden 27 falsch klassifiziert, überwiegend als Stufe~1 oder Stufe~3.

\paragraph{Formale vs.\ inhaltliche Kompetenzsprache.}
\enquote{Sie erhalten Anleitung} ist kein Lernergebnis, sondern beschreibt eine Lehrmethode.
Der Embedding-Ansatz klassifiziert den Satz fälschlich als K (Konfidenz~0{,}940), da das umgebende Fachvokabular (Planung, Dokumentation, Arbeitsprozess) semantisch nah an echten Kompetenzformulierungen liegt.
Der regelbasierte Ansatz erkennt korrekt, dass kein aktives Handlungsverb vorliegt.
Dieses Muster betrifft 30 der 89~Nicht-K-Sätze im Test-Set (33{,}7\,\% Falsch-Positiv-Rate bei I+S).

% =============================================================================
\subsection{Zusammenfassung der Ergebnisse}\label{sec:eval-zusammenfassung}

Tab.~\ref{tab:eval-zusammenfassung} fasst die Kernergebnisse der vergleichenden Evaluation zusammen.

\begin{table}[H]
\caption{Zusammenfassung der Evaluationsergebnisse (F1 weighted)\label{tab:eval-zusammenfassung}}
\centering
\begin{tabular}{lrr}
\hline
\textbf{Ansatz} & \textbf{Typ-Klassifikation} & \textbf{Taxonomie-Zuordnung} \\
\hline
Regelbasiert & 0{,}819 & 0{,}581 \\
Embedding (SVM) & 0{,}963 & 0{,}839 \\
Hybrid & 0{,}929 & 0{,}713 \\
\hline
\end{tabular}
\end{table}

Zur statistischen Absicherung der Unterschiede wurde für alle paarweisen Vergleiche ein McNemar-Test~\cite{mcnemar1947note} durchgeführt (Tab.~\ref{tab:mcnemar}).
Der McNemar-Test prüft auf Basis der diskordanten Paare (Sätze, bei denen genau ein Klassifikator korrekt liegt), ob die Unterschiede zwischen zwei Klassifikatoren statistisch signifikant sind.
Bei sechs Vergleichen wird eine Bonferroni-Korrektur angewandt ($\alpha = 0{,}05/6 \approx 0{,}0083$).

\begin{table}[H]
\caption{McNemar-Tests für paarweise Klassifikatorvergleiche. $n_{01}$/$n_{10}$: diskordante Paare (nur A bzw.\ nur B korrekt). Alle Vergleiche sind nach Bonferroni-Korrektur hochsignifikant ($p < 0{,}001$).\label{tab:mcnemar}}
\centering
\small
\begin{tabular}{llrrrl}
\hline
\textbf{Vergleich} & \textbf{Dimension} & $n_{01}$ & $n_{10}$ & $\chi^2$ & $p$ \\
\hline
Regel vs.\ Embedding & Typ & 11 & 235 & 202{,}15 & $< 0{,}001$ \\
Regel vs.\ Embedding & Taxonomie & 40 & 619 & 506{,}96 & $< 0{,}001$ \\
Regel vs.\ Hybrid & Typ & 8 & 198 & 173{,}40 & $< 0{,}001$ \\
Regel vs.\ Hybrid & Taxonomie & 0 & 358 & 356{,}00 & $< 0{,}001$ \\
Embedding vs.\ Hybrid & Typ & 37 & 3 & 27{,}23 & $< 0{,}001$ \\
Embedding vs.\ Hybrid & Taxonomie & 261 & 40 & 160{,}80 & $< 0{,}001$ \\
\hline
\end{tabular}
\end{table}

Alle sechs Vergleiche sind hochsignifikant ($p < 0{,}001$).
Die diskordanten Paare bestätigen die Rangfolge: Bei Typ-Klassifikation liegt der Embedding-Ansatz in 235~Fällen richtig, in denen der regelbasierte falsch liegt, aber nur in 11~Fällen umgekehrt.
Auch der Unterschied zwischen Embedding und Hybrid ist signifikant ($\chi^2 = 27{,}23$), mit 37~Fällen zugunsten des Embedding-Ansatzes gegenüber 3~zugunsten des Hybrids.

Der \textbf{Embedding-Ansatz} erzielt in beiden Dimensionen die besten Ergebnisse.
Bei der Typ-Klassifikation übertrifft er den regelbasierten Ansatz um 14{,}4~Prozentpunkte, bei der Taxonomie-Zuordnung um 25{,}8~Prozentpunkte.
Der Vorteil ist besonders ausgeprägt bei Sätzen mit unbekannten Verben ($F_1 = 0{,}996$ vs.\ $0{,}683$) und bei Nicht-Kompetenz-Sätzen ($F_1 = 0{,}740$ vs.\ $0{,}279$).

Der \textbf{hybride Ansatz} liegt in beiden Dimensionen zwischen den anderen beiden, erreicht aber nicht die Qualität des reinen Embedding-Ansatzes.
Seine Stärke zeigt sich bei Standardfällen ($F_1 = 1{,}000$) und bei unbekannten Verben ($F_1 = 0{,}976$), wo er die Abdeckung des regelbasierten Ansatzes durch Embedding-Fallback deutlich verbessert.
Die Schwäche liegt im Cascading: Wenn der regelbasierte Ansatz einen Satz fälschlich als K klassifiziert, wird das Embedding nicht mehr konsultiert.

Der \textbf{regelbasierte Ansatz} bleibt bei Standardfällen konkurrenzfähig ($F_1 = 0{,}994$), versagt aber systematisch bei Sätzen außerhalb seines Verblisten-Wissens.
Seine Abdeckung von nur 48{,}2\,\% bei der Taxonomie-Zuordnung zeigt die grundsätzliche Limitation statischer Verblisten.

Diese Ergebnisse werden in Kap.~\ref{sec:diskussion} im Hinblick auf die Forschungsfrage, den Kosten-Nutzen-Aspekt des hybriden Ansatzes und mögliche Verbesserungen diskutiert.
